\section{Fisher's Grouping Corrections and the Conditions on Moments}
\label{sec:p15-fisher-grouping}

The grouping arguments in Fisher's original 1922 paper connect his early
criterion of population consistency to his later likelihood analysis of
discontinuous observations. Their subject is the information actually
recorded: a class, rather than an exact real measurement. Fisher first
calculates how midpoint moments differ from population moments, then asks
whether correcting those moments is as efficient as using the class
probabilities themselves \cite[pp.~317--321,356--363]{Fisher1922Foundations}.
This continuation of \cref{sec:p15-fisher-estimation} separates three
regimes that must not be interchanged: a fixed grid, an average over grid
phase, and a small-width expansion with its remainder. The exact results
below are independently derived; Fisher's historical numerical comparison
and his asymptotic interpretations are identified as such.

\subsection{What is lost when a measurement is replaced by a midpoint?}

\begin{hbdefinition}{def:p15-fg-grid}{Fixed midpoint grouping}
Let \(X\) have density \(f\), let the physical class width be \(h>0\), and
let \(\vartheta\in[0,2\pi)\) be the grid phase. The midpoints and classes are
\[
 c_j=h\left(j+\frac{\vartheta}{2\pi}\right),\qquad
 B_j=[c_j-h/2,c_j+h/2),\qquad j\in\mathbb Z.
\]
The recorded midpoint \(Y_\vartheta\) is \(c_j\) when \(X\in B_j\).
Write \(G_k(\vartheta)=\mathbb E[Y_\vartheta^k]\) and
\(M_k=\mathbb E[X^k]\), when the absolute moments exist. Bars denote
integration over a uniform phase, not a value at a fixed phase:
\(\overline G_k=(2\pi)^{-1}\int_0^{2\pi}G_k(\vartheta)d\vartheta\).
\end{hbdefinition}

The deterministic error \(Y_\vartheta-X\) is bounded by \(h/2\), but it
depends on \(X\). It is not an independent uniform measurement error.
That uniform error is an exact representation only after phase averaging.
This is the content of Fisher's aperiodic term \(A_0\), not of his complete
fixed-grid Fourier series.

\begin{hbtheorem}{thm:p15-fg-phase}{proved=annex:p15.fg.phase}{Exact phase average and Sheppard inversion}
If \(\mathbb E|X|^k<\infty\), then \(G_k\) is integrable in phase and
\(\overline G_k=\mathbb E[(X+U)^k]\), where \(U\) is independent of
\(X\) and uniform on \([-h/2,h/2]\). In particular,
\[
\begin{aligned}
 \overline G_1&=M_1,&
 \overline G_2&=M_2+h^2/12,\\
 \overline G_3&=M_3+h^2M_1/4,&
 \overline G_4&=M_4+h^2M_2/2+h^4/80.
\end{aligned}
\]
Consequently the inverse corrections are
\[
\begin{aligned}
 M_1&=\overline G_1,& M_2&=\overline G_2-h^2/12,\\
 M_3&=\overline G_3-h^2\overline G_1/4,&
 M_4&=\overline G_4-h^2\overline G_2/2+7h^4/240.
\end{aligned}
\]
All moments in these equations are raw moments about the same fixed origin.
\end{hbtheorem}

Fisher obtains these first four terms by integrating over every possible
position of the classes. Omitting the periodic part gives the usual
Sheppard correction; retaining it exposes the approximation
\cite[pp.~318--319]{Fisher1922Foundations}. Applying the displayed inverse
to \(G_k(\vartheta)\) at a particular phase needs a bound on that part.
Centering at a phase-specific estimated mean is another nonlinear operation;
its square cannot be replaced by the square of its phase average.

\begin{hbexample}{def:p15-fg-fixed-failure}{A correction that fails on a fixed grid}
Let \(X\) be uniform on \([-h/2,h/2]\) and use the grid with midpoint zero.
Every recorded value is zero, although the actual variance is \(h^2/12\).
Subtracting \(h^2/12\) from the recorded second moment gives a negative
answer. The phase-average identity remains true. The failure is the
replacement of a fixed-grid functional by its average, not an anomalous
population or a failure of the algebraic inversion.
\end{hbexample}

\subsection{Fisher's periodic term, rather than a rule of thumb}

\begin{hbtheorem}{thm:p15-fg-fourier}{proved=annex:p15.fg.fourier}{Exact normal-grid coefficients and a tail bound}
Suppose \(X\) is normal with mean zero and standard deviation \(\sigma>0\).
For any nonnegative integer \(k\), put
\[
 H(t)=e^{-\sigma^2t^2/2}\frac{2\sin(ht/2)}{t},\qquad
 \omega_s=\frac{2\pi s}{h},\qquad
 C_{k,s}=\frac{i^k}{h}H^{(k)}(\omega_s),
\]
with the continuous value \(H(0)=h\). Then
\[
 G_k(\vartheta)=C_{k,0}
       +2\operatorname{Re}\sum_{s=1}^{\infty}C_{k,s}e^{is\vartheta}
\]
converges absolutely and uniformly in phase. Here \(C_{k,0}=\overline G_k\).
For the first moment this says, with the direction ``recorded minus true,''
\[
 G_1(\vartheta)=\frac h\pi
  \sum_{s=1}^{\infty}\frac{(-1)^{s+1}}s
  e^{-2\pi^2\sigma^2s^2/h^2}\sin(s\vartheta).
\]
Writing \(c=2\pi^2\sigma^2/h^2\), the tail after harmonics \(1,\ldots,m\)
has absolute value at most
\[
 \frac h\pi\,
 \frac{e^{-c(m+1)^2}}{(m+1)(1-e^{-c(2m+3)})}.
\]
The derivative formula supplies the complete coefficients of every higher
moment, not just a leading harmonic.
\end{hbtheorem}

For example, the cosine coefficient of \(G_2\) is
\[
 2\operatorname{Re}C_{2,s}
 =4(-1)^s\left(\sigma^2+\frac{h^2}{4\pi^2s^2}\right)e^{-cs^2}.
\]
The third and fourth coefficients are obtained by taking exactly three
and four derivatives of the same explicitly given \(H\); no endpoint
term is suppressed by that expression. Fisher uses \(h=2\pi\epsilon\)
for these coefficients. His displayed mean ``error'' has the opposite
direction to the recorded-minus-true convention above; a correction to
subtract and a bias to add must not be confused
\cite[pp.~319--320]{Fisher1922Foundations}.

At \(h=\sigma\), the principal first-moment amplitude is
\(\sigma e^{-2\pi^2}/\pi\). Equating it to one tenth of the ungrouped
mean's standard deviation \(\sigma/\sqrt n\) gives
\(n=\pi^2e^{4\pi^2}/100\). Fisher similarly compares leading terms for
moments two through four, obtaining respectively
\(e^{4\pi^2}/800\), \(e^{4\pi^2}/(960\pi^2)\), and
\(3e^{4\pi^2}/(3200\pi^4)\). These are comparisons of a specified
normal-model deterministic harmonic with a sampling standard deviation,
not measured sample requirements or uniform error guarantees. The later
three also drop smaller polynomial terms in their coefficient amplitudes.
The exact coefficient formula and tail above provide the terms needed to
replace these approximations. The source's warning about finite support
endpoints matters: smooth Gaussian decay is doing real work here, and
these exponential bounds do not extend to arbitrary bounded curves.

\subsection{A finite moment is not a finite precision calculation}

Fisher's method-of-moments discussion requires four finite moments to fit
the four-parameter Pearson system, then eight for a finite ordinary
sampling-variance calculation of the fourth sample moment
\cite[p.~321]{Fisher1922Foundations}. These are different conditions.
Neither identifies the curve family merely from four numbers, nor proves
that a fitted parameter map is continuous, unique or efficient.

\begin{hbtheorem}{thm:p15-fg-moments}{proved=annex:p15.fg.moments}{Existence, precision and a heavy-tail failure}
For independent repetitions of the same fixed-grid experiment, the sample
average of \(Y_\vartheta^k\) converges in probability to
\(G_k(\vartheta)\) whenever \(\mathbb E|X|^k<\infty\). Its variance is
\(\operatorname{Var}(Y_\vartheta^k)/n\) if the \(2k\)-th absolute
moment is finite. Finiteness of \(\mathbb E|Y_\vartheta|^r\) and
\(\mathbb E|X|^r\) is equivalent for every \(r>0\).
For the density \(5x^{-6}\) on \([1,\infty)\), the fourth moment is
finite but the eighth is infinite. The fourth sample moment is consistent
in the stated sense but has infinite variance for every finite sample size.
\end{hbtheorem}

The infinite variance is not a proof that no useful estimator exists.
Fisher makes that distinction through his subsequent Cauchy example,
whose actual mean/median contrast and full proof are given in
\cref{sec:p15-fisher-estimation}. Here the Pareto counterexample instead
isolates the precise fourth-versus-eighth moment issue. Even finite
eighth moments alone do not justify an arbitrary nonlinear parameter fit
or a normal approximation.

If the polynomial corrections are applied to sample midpoint moments on
a fixed grid, the preceding consistency result makes their limiting
target the corresponding polynomial in \(G_k(\vartheta)\), not automatically
\(M_k\). Their linear combinations converge because each component
does, by the triangle inequality. Increasing the sample size removes
sampling fluctuation but does not remove the omitted periodic population
bias. That is why a small correction and an exact population calibration
are different claims.

\subsection{Use the probability of the recorded class}

\begin{hbdefinition}{def:p15-fg-cells}{Normalized cell likelihood and discrepancy}
For a fixed partition, let \(p_j(\eta)\) be its cell probabilities, with
\(\sum_jp_j(\eta)=1\). In \(N>0\) independent recorded observations let
\(n_j\) be the cell counts and \(m_j=Np_j(\eta)\). On positive model
cells define
\[
 D(\eta)=\sum_jn_j\log\frac{n_j}{m_j},\qquad
 Q(\eta)=\sum_j\frac{(n_j-m_j)^2}{m_j},
\]
using \(0\log0=0\). The quantity \(D\) is Fisher's \(L\) on p.~357:
negative log-likelihood plus a data-only constant. It is not the
positive log-likelihood used earlier in the estimation argument.
\end{hbdefinition}

\begin{hbtheorem}{thm:p15-fg-deviance}{proved=annex:p15.fg.deviance}{Exact objective and a controlled Pearson approximation}
For finitely many positive differentiable cells, minimizing \(D\) is
equivalent to maximizing the multinomial likelihood. Its score is
\(-\sum_j(n_j/m_j)\partial_\eta m_j\); the derivative of \(Q\) is
\(-\sum_j(n_j^2/m_j^2)\partial_\eta m_j\).
If \(\delta_j=(n_j-m_j)/m_j\) satisfies
\(\max_j|\delta_j|\leq\rho<1\), then
\[
 |2D-Q|\leq\frac{\rho}{3(1-\rho)^2}Q.
\]
For countably many cells the likelihood equivalence holds whenever the
displayed sums exist; derivative exchanges require additional summability.
\end{hbtheorem}

Thus minimizing Pearson's discrepancy is not generally an exact
likelihood operation. Fisher's fine-grouping objection is that squared
counts weight coincidences differently from the likelihood score. His
series gives \(2D\) and \(Q\) the same leading quadratic term, while
retaining the cubic correction \(-\sum_j(n_j-m_j)^3/(3m_j^2)\)
\cite[pp.~357--358]{Fisher1922Foundations}. Empty cells have
\(\delta_j=-1\), and sparse cells may have large positive relative
errors, so the theorem's small-relative-error condition cannot simply be
assumed. Closeness of two objective values does not by itself bound the
distance between their minimizers. Nor does this deterministic Taylor
bound supply a sampling chi-square distribution, degrees-of-freedom
adjustment, or goodness-of-fit calibration. Fisher expressly leaves the
distribution of his \(L\) in other cases for separate investigation.

\subsection{The actual grouped-normal illustration}

\begin{hbdefinition}{def:p15-fg-normal}{The normal class experiment}
For fixed endpoints \(b_j<b_{j+1}\), a normal law with location \(\mu\)
and scale \(\sigma\) gives the exact probability
\[
 p_j=\Phi\left(\frac{b_{j+1}-\mu}{\sigma}\right)
     -\Phi\left(\frac{b_j-\mu}{\sigma}\right).
\]
Writing \(z_j=(b_j-\mu)/\sigma\), differentiation gives
\[
 \partial_\mu p_j=\frac{\phi(z_j)-\phi(z_{j+1})}{\sigma},\qquad
 \partial_\sigma p_j=
 \frac{z_j\phi(z_j)-z_{j+1}\phi(z_{j+1})}{\sigma}.
\]
Here \(\phi\) and \(\Phi\) are the standard normal density and
distribution function.
\end{hbdefinition}
These derivatives follow directly from endpoint
differentiation, not from an unlicensed interchange with a changing
integration domain. The exact likelihood equations use these probabilities,
not a density evaluated at the midpoint. Their finite calculation and
information interpretation are proved in \cref{annex:p15.fg.information}.

Fisher takes a table from K.~Smith's investigation: 300 right-ascension
errors, with positive and negative errors pooled. The absolute-error
classes, in units of \(0.1\) arcsecond as in the source table, and
reported counts are
\[
\begin{array}{c|rrrrrrrrr}
 |x|\text{ class}&0\!:\!1&1\!:\!2&2\!:\!3&3\!:\!4&4\!:\!5&5\!:\!6&6\!:\!7&7\!:\!8&8\!:\!9\\
 n_j&114&84&53&24&14&6&3&1&1.
\end{array}
\]
For the zero-centered model a pooled class probability is
\(2[\Phi((j+1)/\sigma)-\Phi(j/\sigma)]\). The tail class
\([9,\infty)\) must remain in the normalized model even if its observed
count is zero. Renormalizing the nine displayed probabilities instead
would posit observation conditional on \(|X|<9\), a different experiment.
The source does not document such a truncation design.

The midpoint sum is \(\sum_jn_j(j+1/2)^2=1563\), so its uncorrected
second moment is \(1563/300=5.21\); subtracting \(1/12\) gives
\(769/150\). Their square roots give Fisher's reported scale values
\(2.282542\) and \(2.264214\), respectively. He reports a
minimum-\(\chi^2\) scale \(2.355860\), and interpolates the likelihood
score to \(1/\widehat\sigma=0.441624\), hence
\(\widehat\sigma=2.26437\). The score table uses reciprocals
\(0.43,0.44,0.45,0.46\), with reported scores
\(15.135,2.149,-11.098,-24.605\)
\cite[pp.~360--361]{Fisher1922Foundations}.
The first two midpoint calculations are exact finite arithmetic; the
likelihood and minimum-\(\chi^2\) values here are Fisher's reported
historical numerical calculations, not independently certified global
optimizers. Their close agreement in this example is not a universal
comparison theorem. The separate Smith paper's collection, exclusions
and sampling design have not been qualified by this table alone.

\subsection{Two losses and a qualified expansion}

\begin{hbtheorem}{thm:p15-fg-information}{proved=annex:p15.fg.information}{Grouping loses score information}
For a normal location-scale law and any fixed interval partition, let
\(S(X)\) be either scalar score, with the other parameter fixed. The
recorded-class score is \(\mathbb E[S(X)\mid\text{class}]\), and
\[
 I_{\rm full}-I_{\rm class}
   =\mathbb E\operatorname{Var}(S(X)\mid\text{class})\geq0.
\]
For both parameters simultaneously the corresponding difference of score
second-moment matrices is positive semidefinite. This is an exact statement
about the specified experiment, not an asymptotic distribution theorem.
\end{hbtheorem}

Fisher's later small-width calculation first averages the normal density
across a class. To make its aperiodic regime explicit, put
\(a=h/\sigma\), \(z=(y-\mu)/\sigma\), and
\[
 v_a(z)=\int_{-1/2}^{1/2}\phi(z+at)\,dt.
\]
This is the exact density of \((X+U-\mu)/\sigma\) in the independently
phase-randomized experiment. Averaging fixed-grid information over phase
gives the same integrals. It does not assert that every fixed phase gives
those values.

\begin{hbtheorem}{thm:p15-fg-expansion}{proved=annex:p15.fg.expansion}{Controlled aperiodic normal information expansion}
Let \(J_\mu\) and \(J_\sigma\) be the location and scale information
of this phase-averaged experiment, multiplied by \(\sigma^2\).
As \(a\to0\),
\[
\begin{aligned}
 J_\mu&=1-\frac{a^2}{12}+\frac{a^4}{144}
              -\frac{a^6}{1728}+\frac{31a^8}{518400}+O(a^{10}),\\
 J_\sigma&=2\left(1-\frac{a^2}{6}+\frac{a^4}{40}
                  -\frac{a^6}{270}+\frac{83a^8}{129600}\right)
                  +O(a^{10}),\\
 J_\mu^{-1}&=1+\frac{a^2}{12}-\frac{a^8}{86400}+O(a^{10}),\\
 2J_\sigma^{-1}&=1+\frac{a^2}{6}+\frac{a^4}{360}
                  -\frac{a^8}{10800}+O(a^{10}).
\end{aligned}
\]
The remainders are bounded by a constant times \(|a|^{10}\) on each
sufficiently small fixed compact width interval. They refer to these
integrals, not to a sampling-normal approximation at arbitrary \(n\).
\end{hbtheorem}

The exact mean variance in the phase-randomized experiment is
\((\sigma^2+h^2/12)/n\). The exact variance of its raw second moment
about the known true mean is
\(2\sigma^4(1+a^2/6+a^4/360)/n\). Subtracting \(h^2/12\) does not
change that variance. Converting a variance estimator to a scale estimator
and interpreting inverse information as its attainable asymptotic variance
would require a separate delta-method/regular-estimator theorem. The
displayed information and moment expressions establish the comparison of
the corresponding variance benchmarks without asserting such a theorem.

At their leading orders grouping changes the location and scale
benchmarks by \(a^2/12\) and \(a^2/6\). The additional difference from
using the first and corrected second moments is much smaller. The scale
term \(a^8/10800\) agrees with Fisher's printed calculation. His printed
location calculation instead has \(a^6/2880\)
\cite[pp.~361--363]{Fisher1922Foundations}. Within the same stated
aperiodic averaging regime, it does not follow from his preceding
polynomial: substituting the displayed grouped moments into its second
location derivative gives \(1/1728\), rather than the printed
\(1/4320\), in the sixth-order information term. The full derivation in
\cref{annex:p15.fg.expansion} gives the corrected eighth-order difference
\(a^8/86400\). This qualification is not a claim about an unaveraged grid
or a license to discard its periodic part.

\begin{proofboundary}{def:p15-fg-fixed-asymptotic}{Where the approximation stops}
For fixed phase, Gaussian moment harmonics are controlled by
\cref{thm:p15-fg-fourier}; information uses different, nonlinear cell-score
integrands. A uniform quantitative bound on their fixed-phase periodic
remainder has not been established by the moment bound. Consequently
neither the aperiodic expansion nor Fisher's ``less than one per cent''
small-bin statement is adopted as a finite-sample, all-phase precision
guarantee. His endpoint warning also excludes an automatic extension to
densities with a positive or infinite boundary ordinate. Existence and
normality of grouped MLEs, identification after grouping, nuisance-parameter
effects, model adequacy and calibrated goodness-of-fit distributions need
their own premises and proofs. A likelihood for recorded classes does not
settle any of these simply by being maximized.
\end{proofboundary}
